% DOT_IMPORT 13 9.1至9.2_校订续译_工程
\egasection{§9\quad 更换射影嵌入}\source{法稿 \texttt{EGA5\_francais.pdf} 物理第53--55页（手写页码50--52）；英译 \texttt{EGAVaiello.pdf} 第36--38页；第十九版法英对读}
\egap{9.1}
对每个整数 $n>0$，令
\[
P(n)=\mathbb P(\operatorname{Sym}^{n}(E)).
\]
有一个显然的浸入 $u_n:P\to P(n)$，因为在 $S$ 的每个仿射开集上，$\mathcal O_P(n)$ 都由其截面生成，并且
\[
p_*\mathcal O_P(n)\simeq\operatorname{Sym}^{n}(E),
\]
其中 $p:P\to S$ 是投影。\footnote{英译脚注52：“无法辨认。”英译把 $E$ 排作 $e$，把 $\mathcal O_P(n)$ 排作 $0_p(n)$，并漏去关系符。法稿第53页可辨 $E$、$\mathcal O_P(n)$ 及 $p_*(\mathcal O_P(n))$，但接下一行处靠近扫描左边界，关系符本身仍不据此强行恢复。正文沿用有数学依据的典范同构：评估满射 $p^*\operatorname{Sym}^{n}(E)\twoheadrightarrow\mathcal O_P(n)$ 给出相对 Veronese 浸入；在 $E$ 平凡化的开集上，其坐标为所有 $n$ 次单项式。}
若 $f:X\to P$ 是非分歧态射（分别地，是浸入），则 $u_n f:X\to P(n)$ 也如此。有时，研究 $X$ 时用 $u_n f$ 代替 $f$ 是有利的，可以避免 $f$ 在某些方面所具有的十分特殊、且有时令人棘手的行为。
（这种特殊情形的一个例子是8.12(b)中所指出的〔底本原有“sic”标记〕情形：$Y^{\mathrm{sing}}\to P^\vee$ 的像具有维数 $r-1$，但产生的域扩张不可分。另一个例子由 $\mathbb P^3$ 中的二次曲面给出：其所有奇异超平面截面都几何可约，尽管 $X$ 本身几何不可约。）\transnote{这里按原文保留二次曲面这一例子的概括措辞，未扩大为任意特征、任意退化二次曲面的无条件断言。法稿第53页先写“当 $r=2m+1$ 时，维数为 $2m$ 的二次超曲面”，划去后手写为“在 $P^3$ 中”，正文采用后改。这里“二次曲面”指所述光滑二次曲面例子；几何不可约与几何整、约化不能混同。}

\medskip
\egas{9.2}{命题（原稿陈述，所引条件待核）}
设 $S=\operatorname{Spec}(k)$，$X$ 在 $k$ 上光滑，且 $f:X\to P$ 非分歧。取 $n\ge2$，并考虑 $f_n=u_n f$。原稿断言：$f_n:X\to P(n)$ 满足8.8中的等价条件；特别地，在超平面参数空间 $P(n)^\vee$ 中，除去一个余维至少为2的集合后，对于其中的点 $\xi$，相应的超平面截面 $Y_\xi$ 或者光滑，或者只有有限个非光滑点，且这些点都是几何普通奇点。若 $f$ 是浸入，则这样的奇点至多只有一个，并且它在 $k(\xi)$ 上有理。\transnote{法稿第53页在参数空间 $P(n)$ 上清楚有对偶记号；英译漏掉该记号。本轮据法稿恢复 $\xi\in P(n)^\vee$，也与全篇的超平面参数约定一致。此为参数空间记号订正，并非补全8.8条件。底本此处未给出完整证明；下文仅有关于8.8(ii bis)的核验提示。}

\medskip
\noindent\textbf{附记（N.B.）.}
应当以一种不排除 $f$ 并非浸入这一情形的方式来陈述8.8。在8.8的条件(ii bis)下，核验本质上是平凡的。\transnote{“核验本质上是平凡的”为原稿评语，不是本次独立审校的结论。现有底本缺失8.8的标题和完整假设，虽残存含(ii bis)的条件列表，仍无法仅凭这句完成核验，也不能据此声称已恢复缺文。}
无疑，应当在9.1中明确说明：$X$ 相对于 $f_n$ 的超平面截面，无非是用次数为 $n$ 的超曲面代替超平面所得的 $X$ 的“截面”。

% DOT_IMPORT 14 9.3_校订续译_工程

\egas{9.3}{命题（加明固定分支条件的译校表述）}
设 $X$ 为本篇所讨论的有限型 $k$-预概形，几何不可约，且 $\dim f(X)\ge2$；取 $x\in X(k)$ 及 $n\ge2$。对下列三个结论，另加条件
\[
\dim f^{-1}(f(x))\leq \dim X-2.\tag{9.3*}
\]
特别地，当 $f$ 拟有限时，此条件成立。\cmt{D19-R01}{法稿第54页及英译第37页均未写条件（9.3*）；9.2中的非分歧假设也没有被声明为9.3的共同假设。原稿在仅有“$X$ 几何不可约、$\dim f(X)\ge2$、$x\in X(k)$、$n\ge2$”的条件下即断言下面(a)--(c)。这一般不成立：取 $f:\operatorname{Bl}_p\mathbf P^2\to\mathbf P^2$，$x$ 为例外曲线上一点，则经过 $f(x)=p$ 的曲线之总变换都有例外曲线这一固定分支；泛成员还有严格变换，因而可约。正文补入一个明确的充分条件，排除固定余维1分支；它不是原稿的隐含假设，也不声称为必要条件。原证明的全部可读文字仍在下文保留，补充论证及反例见核查报告。}
\begin{enumerate}
\item[(a)] 考虑 $X$ 相对于 $f_n=u_n f$ 的、经过 $x$ 的超平面截面的线性族。其泛成员所定义的 $Y^{(n)}_\xi$ 几何不可约。
\item[(b)] 设 $L$ 是 $P$ 中经过 $f(x)$ 而不包含 $f(X)$ 的线性子簇。考虑 $X$ 相对于 $f_n$ 的如下超平面截面线性族：它由 $P(n)$ 中“在 $x$ 处与 $L$ 相切”的超平面定义，即由 $P$ 上的 $n$ 次齐次式定义，这些式子限制到 $L$ 后在 $f(x)$ 处至少二阶为零。其泛成员 $Y^{(n)}_\xi$ 几何不可约。\transnote{法稿在括号中仍以 $x$ 简写 $f(x)$；正文写出该点真正所属的空间。条件是限制到 $L$ 后在此点二阶消失，不是属于整个 $L$ 的理想平方。}
\item[(c)] 设 $X$ 在 $x$ 处光滑，并且 $n\ge3$，\emph{或} $f(X)$ 不包含于一个在 $k$ 上定义的平面中。考虑 $X$ 相对于 $f_n$ 的、“在 $x$ 处与 $X$ 相切”的超平面截面 $Y^{(n)}_\xi$ 族。其泛成员几何不可约。\cmt{D19-R02}{英译第37页以“where”连接后半句，旧译因而留疑。本轮直接对读法稿第54页，$n\ge3$ 后的手写插入明确以“ou”（或）起首。正文恢复该析取关系；$n\ge2$ 的共同条件仍保留。这个字迹恢复并不消除原稿漏写固定分支条件的问题。}
\end{enumerate}

\noindent\textbf{原稿的证明文字及所需限定。}
用上一节末尾的判据，证明本质上是平凡的。\transnote{这是作者对原证明的评语。上一节有关一般线性系的稠密性推理原有缺条件，本轮已经订正；此处只有在排除固定余维1分支后才能这样应用。补充论证与作者原文分别标明。}
在 $P$ 中取包含 $f(x)$ 的仿射坐标开集，并在其逆像上考虑 $X$。\transnote{法稿第54页清楚写 $f(x)$；英译误排为 $f(X)$。正文恢复法稿的小写 $x$，并非要求整个射影像包含在一个仿射开集中。}
在(a)中，问题归结为：在坐标 $T_1,T_2,\ldots,T_r$ 中找到三个次数不超过2的多项式 $P,Q,R$，使得
\[
\frac{Q(t)}{P(t)},\qquad \frac{R(t)}{P(t)}
\]
在 $K$ 中于 $k$ 上代数无关。这里 $K$ 为 $X$ 的函数域，$t=(t_1,\ldots,t_r)$ 为由 $T_i$ 定义的 $K$ 中元素系。\transnote{如 $X$ 未约化，函数域按其不可约约化概形理解。坐标以 $f(x)$ 为原点。多项式 $P$ 与射影空间 $P$ 同名，沿用原稿。}
在(b)中，还要求 $P,Q,R$ 限制到 $L$ 后在 $f(x)$ 处至少二阶消失；此外可设 $L$ 由方程
\[
T_1=\cdots=T_s=0
\]
定义。\transnote{原稿这里简写为“在 $L$ 上至少二阶消失”。按(b)的明确定义应理解为限制到 $L$ 后在指定点二阶消失；若误解为属于 $I_L^2$，下文 $P=T_1$ 便不满足要求。}
最后，(c)中的情形相同，只是 $L$ 为 $X$ 在 $x$ 处切空间的像；必要时允许取次数为3的 $P,Q,R$，也就是说稍稍放宽次数。

假设 $\dim f(X)\ge2$ 意味着 $k(t_1,\ldots,t_r)$ 在 $k$ 上的超越次数至少为2。因此可找到两个代数无关的坐标函数，不妨记为 $t_1,t_2$。\transnote{法稿第55页为小写 $k(t_1,\ldots,t_r)$。英译误作 $K(t_1,\ldots,t_r)$；因 $t_i\in K$，后者只等于 $K$，不能表达像的维数。}
于是在(a)中取
\[
P=T_1,\qquad Q=T_1^2,\qquad R=T_1T_2.
\]
在(b)中也同样取。注意，此时可选取使 $T_1$ 在 $L$ 上为零的 $t_1$，因为 $f(X)\not\subset L$。这意味着存在 $1\le i\le s$ 使 $t_i\ne0$；由于 $t_i$ 在 $x$ 处为零，它不可能是常数，因而不在 $k$ 上代数。\footnote{英译脚注53：“〔英译者〕由译者补入。”该注附在“不在 $k$ 上代数”之后；法稿第55页也可辨“donc $t_s$ n'est pas alg sur $k$”，故本轮只保存英译脚注的身份，不把这句判为仅见于英译的新增内容。}\textsuperscript{,}\transnote{法稿假设为 $f(X)\not\subset L$；英译误作 $f(x)\not\subset L^*$。法稿和英译都在“存在下标 $i$”后混用了 $t_s$，正文统一为所选的 $t_i$。下面作者允许将基域扩为代数闭域，因而这里“不为常数则不代数”的推理有相应依据。}
（附记：可假设 $k$ 代数闭。）

除去 $f(X)$ 包含在 $X$ 于 $x$ 的切空间经 $f$ 所得的像 $L$ 中这一情形外，(c)由(b)推出。\transnote{法稿第55页此处有一段被划去的文字，英译未收：大意为“这又由维数理由要求 $x$ 处切映射为单射，即 $f$ 在 $x$ 处非分歧，从而 $f(X)$ 包含在 $P$ 中一个与 $X$ 同维的线性子簇中”。其扫描右边界缺损、个别连接词难辨；保留其被删除的身份，不把此段当作新增假设或有效证明。仅凭 $f(X)\subset L$ 并不能推出 $f$ 在 $x$ 处非分歧。}
〔若 $\dim X=2$，这一情形确实是例外：一个二次曲面与其相切平面的交迹一般由两条相交直线组成，所以不可约性不成立。〕
但为处理这一情形，在上面明确给出的 $P,Q,R$ 中，可以把 $X$ 换成 $L$ 本身，而在 $L$ 上问题的解是平凡的。\cmt{D19-R03}{原稿把这一步称为“显然”。在本条 $X$ 于 $x$ 光滑的条件下，可补出所需理由：扩域至代数闭域并令 $m=\dim\overline{f(X)}$，切映射的秩在 $x$ 处至多为 $m$；因为其 $(m+1)$ 阶子式在稠密开集上为零，而 $x$ 的光滑不可约邻域约化，故也在 $x$ 处为零。若 $f(X)\subset L$，则 $m\le\dim L\le m$，因而 $\overline{f(X)}=L$。这才保证 $L$ 上的代数无关函数在像的函数域中仍代数无关。此为译者补证；固定分支的排除仍另需（9.3*）。}
（若 $\dim L=2$，取
\[
P=T_r^2,\qquad Q=T_r^3,\qquad R=T_r^2T_{r-1},
\]
注意 $Q/P=T_r$ 与 $R/P=T_{r-1}$ 是 $L$ 上线性无关的线性式，因而代数无关。）
\orig{法稿第55页随后有已划去的一句：“可以在备注中指出，(c)中甚至可以取 $n=2$，除非 $f(X)$ 包含在一个平面中。”它与第54页后来手写加入的析取条件相呼应。}
\noindent\textbf{英译所载的补充计算。}若 $\dim L\ge3$，则可取在 $L$ 上线性无关的 $T_{r-2},T_{r-1},T_r$，并令\transnote{这段二次式计算见英译第38页，在法稿第55页的相应位置未见同一段文字；本轮保留英译内容并注明来源。两个比值为独立坐标比值的平方，即使特征为2也仍代数无关；不应把不可分性误作代数相关性。英译原括号未闭合，本轮仅整理排版。}
\[
P=T_r^2,\qquad Q=T_{r-1}^2,\qquad R=T_{r-2}^2.
\]

\noindent\textbf{译校补充（不属于原稿）。}三个线性族都包含所有在 $f(x)$ 处至少二阶为零的 $n$ 次形式，所以其基点支集恰为 $f^{-1}(f(x))$。条件（9.3*）使泛截面的余维1分支不可能完全落在基点集中。上述多项式比值保证去掉基点后之映射的像至少二维；应用上一段经限定的 Bertini 不可约性判据，得到该开集上的泛截面几何不可约。其在全截面中稠密，故全截面几何不可约。这里排除的是固定分支；没有由此断言泛截面一定几何整或光滑。

% DOT_IMPORT 15 9.4至10.1_校订续译_工程
\subsection*{（9.4）的编排提示与（9.5）的残文}
结合8.18，我们得到一个推论9.4。（〔英译者〕待陈述。）

最后，还须将后者与9.2结合，以便在“优良情形”下得到一个汇总性的定理。

\egas{9.5}{定理残文（尚无完整陈述）}\footnote{英译标题标“第52页或第55页字迹不清”，原脚注54：“请询问 A.G.，我未能理解。〔英译者〕”法稿物理第55页的手写页码为52。}

〔残文〕（当 $X$ 在 $k$ 上光滑、固有且几何整，$f:X\to P$ 非分歧，$X$ 的维数 $\ge2$）……指出在 $X\to P$ 为浸入时的结果〔此处一个词仍难辨〕（$Y_\xi$ 的奇异点簇）。\transnote{法稿第55页可辨手写“et géom. intègre”、大写 $X\to P$、“immersion”及末尾“variété des pts singuliers de $Y_\xi$”；据此恢复旧英译“$x\to p$”及若干缺口。中间手写词不能可靠确定，且整段仍只是待编定理的提示，未给出完整结论。本轮不据9.2、8.18或10.2反向拼造9.5。}
